\section{Institutional Background}\label{sec:institutional}

The reform that anchors this paper has a simple top-down structure. A generic legal opinion at the São Paulo state Attorney General's office (PGE-SP) reinterpreted a federal procurement-rule threshold in late 2017; the state's centralized procurement platform (BEC-SP) operationalized the new interpretation across all product groups; and the state audit court (TCE-SP) ratified it five months after the empirical cutoff of this paper. The rest of this section explains why the interaction of this transversal reinterpretation with the typical structure of medical-supplies procurement notices mechanically unlocked Group~65 from the SME-only default, providing the natural experiment exploited below. Table~\ref{tab:institutional_timeline} reconstructs the chronology.

Brazil's federal SME statute (Complementary Law 123/2006, amended for procurement by \LCSetAsideAct{} and for revenue brackets by \LCRevenueAct{}) defines micro and small enterprises by annual gross revenue ceilings of R\$\MERevCeilingBRL{} (US\$\MERevCeilingUSD{}) and R\$\EPPRevCeilingBRL{} (US\$\EPPRevCeilingUSD{} million), respectively, and mandates exclusive SME tenders for items valued at R\$\setAsideThreshBRL{} or less. \LCSetAsideAct{} converted the SME-only rule from optional (``may'') to mandatory (``must''), with departures requiring a costly justification to the audit courts; state-level adherence to SME-only tendering for eligible items rose from roughly \adherencePreLC\% to nearly \adherencePostLC\% in the years that followed.\footnote{Adherence figures and Group-65 sub-rates are author tallies from the BEC microdata extract used in the structural sample (Section~\ref{sec:data}; replication numbers in Table~\ref{tab:sample_flow}), restricted to items satisfying the federal R\$\setAsideThreshBRL{} threshold. The R\$\EPPRevCeilingUSD M EPP ceiling raised by \LCRevenueAct{} has no direct bite on the structural sample, where SME status is identified from the BEC \emph{enquadramento} flag and the Receita Federal \emph{porte} register at the firm-auction level; the bracket change is reported here for institutional completeness.}

The state of São Paulo has operated the Bolsa Eletrônica de Compras (BEC) since September 2000 (Decree \decreeBECOne), and the Pregão electronic reverse-auction format that anchors this paper's analytical sample was regulated within BEC by Decree \decreePregao. BEC centralizes Pregão for standardized-goods procurement by state bureaus and affiliated units; the structural sample covers \itemTransactionsCount{}~million item-level transactions across all product groups over the \nDataMonths-month estimation window, of which Group~65 accounts for roughly \groupSixtyfivePctTx{}~percent. Its catalog is organized in three tiers---group, class, and item---so that ``medical, dental, and hospital equipment and supplies'' form Group~65, within which class~6531 identifies medications subject to price ceilings set by the Câmara de Regulação do Mercado de Medicamentos (CMED).\footnote{Brazilian institutional terms used throughout the paper (Pregão, Convite, BEC, CADMAT, CMED, PGE-SP, TCE-SP, Sefaz-SP) are defined in Appendix~\ref{app:glossary}. Platform institutional description at \url{www.bec.sp.gov.br/BECSP/Quem_Somos/Quem_Somos_BEC_novo.aspx} (retrieved April 2026); the \itemTransactionsCount-million item-level count and the \purchaseOffersCount{} purchase-offer count are author tallies from the same BEC extract used in the structural analysis.}

Group 65 was a singular exception to the 2014 rule, but for a transversal legal reason rather than a sector-specific one. Until late 2017, the prevailing interpretation of art.~48,~I of \LCSmeStatute{} was that the R\$\setAsideThreshBRL{} threshold for SME-only tenders applied to the \emph{total estimated value of the procurement notice}, not to each individual item or lot---a reading sustained by TCE-SP in eTC-\eTCRestrictive{} (23 September 2015, by tie-breaking vote). Group-65 procurement notices systematically combine small individual items (medications, supplies, dressings) with large notice totals reflecting bulk public-sector purchasing, so the totalizing reading mechanically barred Group~65 from the SME-only default---not because medical supplies were classified as ``strategic,'' but because their typical procurement structure was exactly what the threshold-on-total rule excluded.

On 11 December 2017, PGE-SP issued \textbf{Parecer Sub-G Cons.\ n\textordmasculine~\pareceNum}, holding that the R\$\setAsideThreshBRL{} threshold applies item-by-item---allowing mixed-exclusivity procurement notices in which items below R\$\setAsideThreshBRL{} are SME-only and items above remain open---and that the more generous federal regime introduced by \LCSetAsideAct{} supersedes the previously applicable São Paulo state acts (Lei Estadual \leiEstadualSp{} and Decree \decreeStatePrev) that had embedded the restrictive interpretation in state-level practice.\footnote{Parecer Sub-G Cons.\ n\textordmasculine~151/2017, Subprocuradoria-Geral de Consultoria da PGE-SP, ementa: \emph{``Licitações exclusivas para ME e EPP (artigo 48, I, LC nº 123/2006) devem ser realizadas para itens ou lotes cujo valor estimado, individualmente, não ultrapasse o limite legal de R\$80.000,00, \textbf{independentemente do valor total da Oferta de Compra}. \ldots\ Aplicabilidade, na íntegra, da normativa federal.''}}\footnote{Disclosed at the BEC/SP legislation portal (\url{https://www.bec.sp.gov.br/becsp/Legislacao/UI_Selecao.aspx}, \emph{Pareceres} section, retrieved April 2026); the Boletim CEPGE indexes only Pareceres da Procuradoria Administrativa, so the BEC portal is the canonical channel for Sub-G Cons.\ pareceres.}

Operationally, BEC moved before PGE consolidated the legal interpretation. COMUNICADO BEC n\textordmasculine~\comunicadoBECone{} (18 July 2017) introduced the SME-only OC functionality, and COMUNICADO BEC n\textordmasculine~\comunicadoBECtwo{} (20 November 2017) expanded it to mixed-exclusivity Pregão notices combining ample-participation, 25\%-quota, and SME-only items in the same notice---the precise operational mechanism that Parecer \pareceNum{} ratified three weeks later.\footnote{Both Comunicados BEC are disclosed at the BEC/SP legislation portal cited above (\emph{Comunicados} section). TCE-SP itself reversed its restrictive reading on 16 May 2018 in eTC-\eTCExpansive{} (Tribunal Pleno), aligning with PGE-SP's expansive interpretation two and a half months after the empirical cutoff of this paper. Both TCE-SP rulings are accessible at the court's jurisprudence portal (\url{https://www.tce.sp.gov.br/jurisprudencia}).} The empirical cutoff of \policyCutoffMonth{} therefore captures the moment of mass take-up of the new functionality by Group-65 PBUs rather than formal regulatory enablement (which was incremental between July and November 2017); SME-only adherence within Group 65 jumped to \adherenceGsixtyfivePost\% on average in the months following the cutoff, below the rate for other groups largely because class-6531 pharmaceuticals are oligopolistic and often fail the three-eligible-SMEs requirement.

What matters for identification is that the trigger was a generic statutory reinterpretation, and the unlocking of Group-65 set-asides was a mechanical consequence of how medical-supplies notices are structured. PGE-SP's parecer made no substantive claim about medical supplies or any other product group; it pronounced on a general legal question, with timing set by the bureaucratic process inside PGE-SP and by when BEC operationalized the new interpretation. The TCE-SP reversal in May 2018 lies after the cutoff and confirms that the trigger was a top-down legal interpretation: when the audit court itself flipped, it ratified what PGE-SP and BEC had already established for state-level procurement rather than initiating an independent change in the Group-65 environment. The change is therefore exogenous to Group-65 market conditions on the dimensions the design tests. This is the institutional basis for Assumption~\ref{a:setaside}: the cutoff shifts admissibility and equilibrium entry selection, not the underlying production technology of the goods themselves. Three empirical checks corroborate the institutional reading: the reduced-form DiD against the \controlGroupCount{} never-treated product groups (Appendix~\ref{app:did}); the primitive-invariance KS test on non-SME $F_c$ (Table~\ref{tab:v3_primitive_invariance}); and the cross-modality consistency between Convite GPV inversion and Preg\~ao drop-out recovery in the unaffected pharma non-SME Pre stratum (Figure~\ref{fig:v3_cross_modality_uh}).

\begin{table}[!htbp]
\centering
\begin{threeparttable}
\caption{Institutional timeline---federal and São Paulo state acts relevant to the structural sample.}
\label{tab:institutional_timeline}
\small
\setlength{\tabcolsep}{3pt}
\begin{tabular}{@{}l l p{0.40\linewidth}@{}}
\toprule
Date & Act & Content \\
\midrule
Sep~2000  & Decree \decreeBECOne{} (SP)              & BEC platform created; renamed Decree \decreeBECTwo. \\
Jun~2005  & Decree \decreePregao{} (SP)              & Pregão electronic reverse-auction regulated within BEC. \\
Dec~2006  & \LCSmeStatute{} (federal)                & SME statute (\emph{Estatuto Nacional}). \\
2008      & Lei Estadual \leiEstadualSp{} (SP)        & SP state-level SME procurement regulation (less favorable than federal post-2014). \\
2009      & Decree \decreeStatePrev{} (SP)              & Regulates Lei Estadual \leiEstadualSp. \\
Aug~2014  & \LCSetAsideAct{} (federal)                & SME-only default (``may''$\to$``must''); R\$\setAsideThreshBRL{} threshold. \\
Sep~2015  & TCE-SP eTC-\eTCRestrictive{}               & Restrictive reading of art.~48,~I (R\$\setAsideThreshBRL{} on notice total) sustained by tie-breaking vote. \\
Oct~2016  & \LCRevenueAct{} (federal)            & EPP revenue ceiling R\$3.6M$\to$R\$\EPPRevCeilingUSD M (effective Jan~2018). \\
18 Jul 2017 & COMUNICADO BEC n\textordmasculine~\comunicadoBECone{} & SME-only OC functionality enabled in BEC. \\
20 Nov 2017 & COMUNICADO BEC n\textordmasculine~\comunicadoBECtwo{} & Mixed-exclusivity (Ample / 25\% quota / SME-only) item-by-item enabled. \\
11 Dec 2017 & PGE-SP Parecer Sub-G Cons.\ \pareceNum{}      & Item-by-item R\$\setAsideThreshBRL{} reading; federal regime supersedes Lei \leiEstadualSp. \\
\policyCutoffMonth{}  & Mass adoption (empirical cutoff)\tnote{a}    & Group-65 PBU take-up of the new functionality (first wave of adopting PBUs). \\
16 May 2018 & TCE-SP eTC-\eTCExpansive{}                    & Tribunal Pleno reverses, aligns with PGE-SP expansive reading. \\
Apr~2021  & \leiNova{} (federal)            & New procurement statute; SME defaults preserved. \\
\bottomrule
\end{tabular}
\begin{tablenotes}\footnotesize
\item[a] BEC formally enabled the new functionality on 18 July 2017 (Comunicado \comunicadoBECone) and expanded it to item-by-item mixed exclusivity on 20 November 2017 (Comunicado \comunicadoBECtwo). The empirical cutoff in \policyCutoffMonth{} (BEC purchase-order sequence index $\mathrm{data\_oc\_numb} = \dataOcCutoff$) corresponds to the moment of mass take-up by Group-65 PBUs, not formal enablement; the gap reflects the diffusion lag between operational availability and PBU adoption. The TCE-SP reversal in May 2018 (eTC-\eTCExpansive) post-dates the cutoff, confirming that the trigger of the design is the PGE-SP/BEC operationalization sequence, not the audit-court realignment.
\end{tablenotes}
\end{threeparttable}
\end{table}
